\documentclass[10pt,a4paper]{amsart}
\usepackage[margin=19mm]{geometry}
\usepackage{amsmath,amssymb,mathtools}
\usepackage[scr=rsfs,cal=euler]{mathalfa}
\usepackage{xfrac}
\usepackage[unicode,hidelinks,bookmarks=false]{hyperref}
\newcommand{\ie}{i.e.\ }
\newcommand{\cf}{cf.\ }
\newcommand{\resp}{respectively\ }
\newcommand{\Subsection}{\subsection}
\providecommand{\nobreakdash}{\nobreak\hbox{-}}
\setlength{\emergencystretch}{2em}
\multlinegap0pt
\usepackage{amssymb}
\usepackage{multirow}
\usepackage{bm}
\usepackage{float}
\usepackage[shortlabels]{enumitem}
\usepackage{xcolor}
\usepackage{mathtools}
\usepackage{bbm}
\usepackage{booktabs}
\newtheorem{lemma}{Lemma}
\title{A numerical study for tempered time-fractional advection-dispersion equation on graded meshes}
\author{Liangcai Huang, Lin Li, Shujuan L\"u}
\date{Source: arXiv:2602.08325v1 (CC BY 4.0)}
\begin{document}
\maketitle

\noindent\emph{Source and licence.} A numerical study for tempered time-fractional advection-dispersion equation on graded meshes. Version arXiv:2602.08325v1 (CC BY 4.0), distributed by arXiv under the Creative Commons Attribution 4.0 International licence, http://creativecommons.org/licenses/by/4.0/, as recorded in the arXiv licence metadata for this exact version and shown on the abstract page https://arxiv.org/abs/2602.08325v1. Full licence notice: \texttt{licenses/arxiv-2602.08325-CC-BY-4.0.txt}.\par\smallskip

\noindent\textit{Editorial scope.} Counted unit: the article's lemma dnd0 (source0.tex lines 569--574). The article's complete original proof of it is delivered with it (source0.tex lines 575--616, one \texttt{proof} environment of 42 lines). No mathematics is written, repaired or re-derived: the statement and the proof are the article's own lines, in the article's own notation, and no retained line is substituted or re-flowed.  Retained uncounted as the source-local context the proof needs: lines 213--215; lines 238--245; lines 254--259; lines 559--568.  Nothing was omitted from any retained span, so the package carries no omission record.  No retained line contains a citation, so the wrapper carries no bibliography and no bibliography entry.  No retained line contains a figure environment, an image-inclusion command or drawing code, so no image is delivered and the package declares no asset dependency.  Editorial formatting only, in addition: an AMS \texttt{amsart} preamble that restates the article's own declarations and macros used by the retained lines, namely the environments \texttt{lemma} on separate counters, together with the standard packages those lines need.  The retained environments print in this document as Lemma 1 in order of appearance, whereas the article prints the same items with the numbering of its own full document.  The complete native source of the article as distributed in the arXiv e-print for this exact version is bundled unchanged as \texttt{sources/194239/source0.tex}.
\begin{align}
    |t^{-1-\alpha}-\sum_{l=1}^{N_{exp}}\omega_le^{-s_lt}|\leq \varepsilon, \quad \forall t\in [\Delta t,T].\label{SOE}
\end{align}

\begin{equation}\label{lambda}
    \begin{aligned}
        \lambda_{i, n}^1
        &=\frac{e^{-(\lambda+s_i){\tau_{n+1}}/{2}}}{(\lambda+s_i)^2\tau_{n}}(e^{-(\lambda+s_i) \tau_{n}}-1+(\lambda+s_i) \tau_{n}), \\
        \lambda_{i, n}^{2}
        &=\frac{e^{-(\lambda+s_i){\tau_{n+1}}/{2}}}{(\lambda+s_i)^2\tau_{n}}(1-e^{-(\lambda+s_i) \tau_{n}}-e^{-(\lambda+s_i) \tau_{n}}(\lambda+s_i) \tau_{n}),
    \end{aligned}
\end{equation}

\begin{equation}\label{ab}
\begin{aligned}
&a_{j,  n}=\alpha \sum_{i=1}^{N_{\exp }} \omega_i e^{-(\lambda+s_i)(\bar{t}_n-\bar{t}_{n-j})} \lambda_{i,  n-j}^1,  \\
&b_{j,  n}=\alpha \sum_{i=1}^{N_{\exp }} \omega_i e^{-(\lambda+s_i)(\bar{t}_n-\bar{t}_{n-j})} \lambda_{i,  n-j}^2. 
\end{aligned}
\end{equation}

\begin{equation}\label{d0}
\left\{
\begin{aligned}
    &(-\frac{1}{h^2}\cos(h\beta)+\frac{1}{2h}i\sin(h\beta)+(\frac{1}{\tau_{n+1}}+\frac{\tau_{n+1}^{-\alpha}}{2^{1-\alpha}\Gamma(2-\alpha)}+\frac{1}{h^2}) )\hat{\rho}^{n+1}[k] \\
    &=(\frac{1}{h^2}\cos(h\beta)-\frac{1}{2h}i\sin(h\beta)+(\frac{1}{\tau_{n+1}}-\frac{\tau_{n+1}^{-\alpha}}{2^{1-\alpha}\Gamma(2-\alpha)}-\frac{1}{h^2}) )\hat{\rho}^{n}[k] \\
    &\quad +\frac{1}{\Gamma(1-\alpha)}((a_{0, n} +\frac{e^{-\lambda {\tau_{n+1}}/{2}}\tau_{n+1}^{-\alpha}}{2^{-\alpha}(1-\alpha)}-e^{-\lambda {\tau_{n+1}}/{2}} ({\tau_{n+1}}/{2})^{-\alpha})\hat{\rho}^{n}[k]\\
    &\quad +\sum\limits_{l=1}^{n-1} (a_{n-l, n}+b_{n-1-l, n}) \hat{\rho}^{l}[k]+(b_{n-1, n}+{e^{-\lambda \bar{t}_n}}{\bar{t}_n^{-\alpha}}) \hat{\rho}^{0}[k]).
\end{aligned}
\right.
\end{equation}

\begin{lemma}\label{dnd0}
    Suppose that $\hat{\rho}_{n}[k]$ $ (1\le n\le N)$ satisfy (\ref{d0}).  Then 
    \begin{align}
    |\hat{\rho}^{n}[k]|\leq |\hat{\rho}^{0}[k]|,\quad 1\le n\le N.\label{rho1}
    \end{align}
\end{lemma}

\begin{proof}
Use induction on $n$. The case $n=1$ of (\ref{rho1}) is
\begin{align}
    |\hat{\rho}_{1}[k]|=\left|\frac{\frac{1}{h^2}\cos(h\beta)-\frac{1}{2h}i\sin(h\beta)+(\frac{1}{\tau_{1}}-\frac{1-2e^{-\lambda {\tau_{1}}/{2}}}{2^{1-\alpha}\Gamma(2-\alpha){\tau_{1}}^\alpha}-\frac{1}{h^2}) }{-\frac{1}{h^2}\cos(h\beta)+\frac{1}{2h}i\sin(h\beta)+(\frac{1}{\tau_{1}}+\frac{1}{2^{1-\alpha}\Gamma(2-\alpha){\tau_{1}}^\alpha}+\frac{1}{h^2}) }\right||\hat{\rho}^{0}[k]| \leq |\hat{\rho}^{0}[k]|.
\end{align}
Fix $n\in \{1,2,\dots,N-1\}$. Assume that (\ref{rho1}) is valid for \(m=1,2,\dots,n\). Then (\ref{lambda}), (\ref{ab})
and the inductive hypothesis yield
\begin{align}
  a_{0,n}+\sum_{l=1}^{n-1}(a_{l, n}+b_{l, n})+b_{n-1,n}  
  &=\sum_{l=0}^{n-1}\alpha \sum_{i=1}^{N_{exp}}w_i\int_{t_{n-l-1}}^{t_{n-l}}e^{-(\lambda+s_i)(\bar{t}_n-s)}ds.\notag
\end{align}
Combining the SOE approximation \eqref{SOE} and the above expression yields
\begin{align}
  a_{0,n}+\sum_{l=1}^{n-1}(a_{l, n}+b_{l, n})+b_{n-1,n}  
&=\alpha\int_{{\tau_{n+1}}/{2}}^{\bar{t}_n}{e^{-\lambda s}}{s^{-1-\alpha}}ds\notag\\
  &\leq {e^{-\lambda {\tau_{n+1}}/{2}}}{({\tau_{n+1}}/{2})^{-\alpha}}-{e^{-\lambda \bar{t}_n}}{\bar{t}_n^{-\alpha}}.\label{aplusb}
\end{align}
Applying the inequality (\ref{aplusb}) to (\ref{d0}), we get
\begin{align}
    &|-\frac{1}{h^2}\cos(h\beta)+\frac{1}{2h}i\sin(h\beta)+(\frac{1}{\tau_{n+1}}+\frac{\tau_{n+1}^{-\alpha}}{2^{1-\alpha}\Gamma(2-\alpha)}+\frac{1}{h^2}) ||\hat{\rho}^{n+1}[k]|\notag \\
    \leq&|\frac{1}{h^2}\cos(h\beta)-\frac{1}{2h}i\sin(h\beta)+(\frac{1}{\tau_{n+1}}-\frac{\tau_{n+1}^{-\alpha}}{2^{1-\alpha}\Gamma(2-\alpha)}-\frac{1}{h^2}) ||\hat{\rho}^{n}[k]| +\frac{e^{-\lambda {\tau_{n+1}}/{2}}\tau_{n+1}^{-\alpha}}{2^{-\alpha}\Gamma(2-\alpha)}|\hat{\rho}^{0}[k]|.\label{need1}
\end{align}
Define
\begin{align*}
A &= \frac{1-\cos(h\beta)}{h^2},\quad B=\frac{\tau_{n+1}^{-\alpha}}{2^{1-\alpha}\Gamma(2-\alpha)},\quad C=2Be^{-\lambda{\tau_{n+1}}/{2}},\\
a&=-A+\frac{1}{\tau_{n+1}}-B,\quad b=A+\frac{1}{\tau_{n+1}}+B,\quad d=\frac{\sin(h\beta)}{2h}.
\end{align*}
To prove that $|\hat{\rho}^{n+1}[k]|\leq |\hat{\rho}^{0}[k]|$ by \eqref{need1}, it suffices to establish the following inequality:
\begin{align*}
& \sqrt{a^2 + d^2} + C \le \sqrt{b^2 + d^2}.
\end{align*}
Noticing that $d^2 \leq {A}/{2}$ and $(A+B)^2-B^2e^{-\lambda\tau_{n+1}}\geq A^2+2AB$, it is sufficient to verify
\begin{align}
    {Be^{-\lambda\tau_{n+1}}\tau_{n+1}^2}\leq 4({1-B^2e^{-\lambda\tau_{n+1}}\tau_{n+1}^2}).\label{need2}
\end{align}
Using the bound \(1/2<\Gamma(2-\alpha)<1\) and straightforward estimates, inequality (\ref{need2}) can be reduced to the simpler condition
\[
\tau_{n+1}^{2-2\alpha}<1/3,
\]
which is easily satisfied in typical numerical simulations. Consequently, (\ref{need1}) and (\ref{need2}) yield \(|\hat{\rho}^{n+1}[k]|\leq |\hat{\rho}^{0}[k]|\). Thus we
have proved (\ref{rho1}) for $m = n+1$. By the principle of induction, the lemma is proved.
\end{proof}

\end{document}
